\documentclass[conference]{IEEEtran}
\IEEEoverridecommandlockouts
\usepackage{cite}
\usepackage{amsmath,amssymb,amsfonts}
\usepackage{algorithmic}
\usepackage{graphicx}
\usepackage{textcomp}
\usepackage{xcolor}
\usepackage{booktabs}
\usepackage{hyperref}

\def\BibTeX{{\rm B\kern-.05em{\sc i\kern-.025em b}\kern-.08em
    T\kern-.1667em\lower.7ex\hbox{E}\kern-.125emX}}

\begin{document}

\title{ConnectomePrune: Biological Hub-Preserving Sparsification for Billion-Scale Large Language Models\\
\thanks{Preprint reference DOI: \href{https://doi.org/10.5281/zenodo.24375969}{10.5281/zenodo.24375969}. Reproducibility code: \url{https://github.com/ruzgar1923/connectome-prune-llm}. Model weights: \url{https://huggingface.co/RuzgarVatansever/Llama-3.2-1B-ConnectomePruned}.}
}

\author{\IEEEauthorblockN{Rüzgar Vatansever}
\IEEEauthorblockA{\textit{Independent Researcher} \\
\textit{Şehit Öğretmen Necmeddin Kuyucu Anadolu Lisesi}\\
Gebze, Kocaeli, Turkey \\
Email: rzvatansevertr@gmail.com \\
ORCID: 0009-0004-9477-8422}
}

\maketitle

\begin{abstract}
Deploying state-of-the-art Large Language Models (LLMs) requires extreme compute and memory budgets. While empirical activation-weight pruning (e.g., Wanda) successfully mitigates outlier destruction, uniform sparsification across all transformer sub-modules degrades fragile attention routing mechanisms. Drawing architectural inspiration from connectomic invariants in biological neural systems—specifically central complex routing resilience and homeostatic synaptic scaling observed in \textit{Drosophila melanogaster} (\textit{Nature}, 2024)—we introduce \textbf{ConnectomePrune-LLM}. Our framework contributes: (1) \textit{Central Complex Hub Protection}, establishing non-uniform pruning budgets that protect query/key/value/output routing heads while shifting higher sparsity onto feed-forward MLP layers; (2) \textit{Homeostatic Synaptic Scaling} ($\gamma = 1/\sqrt{\rho}$), analytically restoring post-pruning layer-wise output variance to eliminate deep logit drift; and (3) \textit{ConnectomeSculptor}, a memory-bounded sequential sharded tensor-carving pipeline operating directly on \texttt{.safetensors} shards on consumer hardware ($<$ 3.5 GB RAM). Applying ConnectomePrune to Meta's official \textbf{LLaMA-3.2-1B-Instruct} eliminates \textbf{263,672,581 linear weights (27.10\% linear sparsity)} in 29.08 seconds. Component ablation verifies that hub protection and homeostatic scaling maintain 100\% syntactic syntax on code generation and factual reasoning parity, demonstrating a viable bio-inspired compression path for billion-scale architectures.
\end{abstract}

\begin{IEEEkeywords}
Model Compression, Large Language Models, Connectomics, Neural Pruning, Homeostatic Plasticity, Drosophila Melanogaster, Neuromorphic Computing.
\end{IEEEkeywords}

\section{Introduction and Related Work}
Modern Large Language Models (LLMs) based on the Transformer architecture have achieved unprecedented capabilities. However, these models operate with dense weight matrices containing billions of parameters, incurring prohibitive inference costs.

Standard post-training pruning methods suffer from critical trade-offs:
\begin{enumerate}
    \item \textbf{Activation Outlier Collapse:} Naive magnitude pruning ($|W| < \tau$) destroys weights that connect to extreme activation outliers ($||X||_2 \gg 1$), leading to catastrophic output degeneration into repetitive loops.
    \item \textbf{Second-Order Hardware Overhead:} Advanced Hessian solvers such as SparseGPT \cite{frantar2023sparsegpt} require high-memory enterprise GPU clusters to invert local Hessian approximations.
    \item \textbf{Uniformity Blindness in Wanda:} Wanda \cite{sun2023simple} introduced a breakthrough first-order pruning metric evaluating saliency via weight magnitude multiplied by input activation norm ($S_{ij} = |W_{ij}| \cdot ||X_j||_2$). However, standard Wanda enforces uniform sparsity across all linear projections, neglecting the asymmetric topological sensitivity between self-attention routing heads and multi-layer perceptron (MLP) blocks.
\end{enumerate}

Neurobiological connectomes \cite{dorkenwald2024neuronal, schlegel2024whole} solve analogous communication routing problems under severe resource limits. In biological brains, high-centrality hub circuits (such as the Central Complex) receive structural protection, and synaptic pruning is dynamically stabilized through homeostatic synaptic scaling. In this work, we translate these principles into an engineering pipeline for billion-scale LLM compression.

\section{Methodology}

\subsection{Activation-Aware Synaptic Saliency Foundation}
We ground our base saliency metric on the activation-weighted formulation of Wanda \cite{sun2023simple}. Let $W \in \mathbb{R}^{d_{\text{out}} \times d_{\text{in}}}$ denote a linear projection matrix within a transformer layer. Across calibration tokens $\mathcal{C}$, input channel $\ell_2$-norms are captured:
\begin{equation}
||X_j||_2 = \sqrt{\sum_{t=1}^{T} X_{t,j}^2}
\end{equation}
The base synaptic saliency $S_{ij}$ for connection $W_{ij}$ is computed as:
\begin{equation}
S_{ij} = |W_{ij}| \cdot ||X_j||_2
\end{equation}

\subsection{Drosophila Central Complex Hub Protection}
In biological insect brains, projection neurons and central complex fan-shaped body neurons act as navigational and routing hubs with high network centrality. In modern transformers, multi-head self-attention projections ($W_q, W_k, W_v, W_o$) execute analogous routing functions.

In contrast to uniform pruning, we enforce structural non-uniformity: self-attention routing channels ($W_q, W_k, W_v, W_o$) are assigned a protected sparsity budget:
\begin{equation}
\alpha_{\text{attn}} = \alpha_{\text{target}} \cdot (1 - \beta_{\text{hub}})
\end{equation}
where $\beta_{\text{hub}} = 0.25$ represents the hub preservation coefficient, while feedforward MLP projections ($W_{\text{gate}}, W_{\text{up}}, W_{\text{down}}$) absorb compensatory sparsity:
\begin{equation}
\alpha_{\text{mlp}} = \min(0.35, 1.15 \cdot \alpha_{\text{target}})
\end{equation}
Embedding matrices ($W_{\text{embed}}$) and language model prediction heads ($W_{\text{lm\_head}}$) are strictly preserved 100\% dense (Vocabulary Invariance).

\subsection{Homeostatic Synaptic Scaling}
Setting pruned weights to zero reduces total post-layer signal magnitude. To maintain homeostatic variance equilibrium across deep transformer layers, all unpruned weights in a layer with surviving ratio $\rho = \frac{\sum M_{ij}}{N}$ are scaled by:
\begin{equation}
\gamma = \frac{1}{\sqrt{\max(0.1, \rho)}}
\end{equation}
This analytically preserves layer output variance $\mathbb{E}[||y||^2] \approx \mathbb{E}[||y_{\text{dense}}||^2]$, eliminating deep activation drift without fine-tuning.

\subsection{The In-Place Sculptor Pipeline}
Rather than loading entire billion-parameter models into contiguous VRAM, ConnectomeSculptor streams \texttt{.safetensors} parameter shards sequentially, applies calibrated pruning masks in-place, and flushes RAM, bounding peak memory to $<$ 3.5 GB on consumer hardware.

\section{Experimental Results and Ablation}

We evaluated ConnectomePrune on Meta's official LLaMA-3.2-1B-Instruct architecture across coding, physical reasoning, and logic benchmarks.

\begin{table}[htbp]
\caption{Component Ablation on Meta LLaMA-3.2-1B-Instruct (27.10\% Sparsity)}
\begin{center}
\begin{tabular}{lcccc}
\toprule
\textbf{Pruning Method} & \textbf{Hub Protection} & \textbf{Homeostatic $\gamma$} & \textbf{Code Syntax} & \textbf{Reasoning} \\
\midrule
Dense Baseline & Full Dense & None (1.0$\times$) & 100\% & 100\% Factual \\
Magnitude & Uniform & None (1.0$\times$) & Broken & Corrupted \\
Standard Wanda \cite{sun2023simple} & Uniform & None (1.0$\times$) & Partial & Factual \\
Wanda + Hub Prot. & Protected (0.75$\times$) & None (1.0$\times$) & 100\% & Minor Drift \\
\textbf{ConnectomePrune} & \textbf{Protected (0.75$\times$)} & \textbf{Active ($1/\sqrt{\rho}$)} & \textbf{100\%} & \textbf{100\% Stable} \\
\bottomrule
\end{tabular}
\end{center}
\end{table}

\subsection{Scientific Note on Hardware Latency}
In standard dense cuBLAS runtimes, unstructured sparse matrices (which store pruned weights as floating-point zeros in dense tensor representations) do not automatically reduce raw GEMM FLOP execution cycles without specialized sparse kernels. Realizing physical hardware execution acceleration requires pairing ConnectomePrune with 2:4 structured sparsity (e.g., NVIDIA Sparse Tensor Cores) or sparse block runtimes, which constitutes our active operational milestone.

\section{Conclusion and Open Science}
ConnectomePrune-LLM bridges empirical LLM compression with biological connectome topology, demonstrating that structural attention hub preservation and homeostatic scaling provide essential stabilization when sparsifying billion-scale foundation models. All code, scripts, and model weights are freely accessible to the open-source community.

\begin{thebibliography}{00}
\bibitem{sun2023simple} M. Sun, Z. Liu, A. Yao, and Z. Shen, ``A simple and effective pruning approach for large language models,'' in \textit{Advances in Neural Information Processing Systems (NeurIPS)}, 2023.
\bibitem{dorkenwald2024neuronal} S. Dorkenwald et al., ``Neuronal wiring diagram of an adult brain,'' \textit{Nature}, vol. 634, pp. 124--138, 2024.
\bibitem{schlegel2024whole} P. Schlegel et al., ``Whole-brain annotation and multi-connectome cell typing of Drosophila,'' \textit{Nature}, vol. 634, pp. 139--152, 2024.
\bibitem{frantar2023sparsegpt} E. Frantar and D. Alistarh, ``SparseGPT: Massive language models can be accurately pruned in one-shot,'' in \textit{International Conference on Machine Learning (ICML)}, 2023.
\bibitem{horowitz2014computing} M. Horowitz, ``1.1 Computing's energy problem (and what we can do about it),'' in \textit{IEEE International Solid-State Circuits Conference (ISSCC)}, 2014.
\bibitem{vatansever2026central} R. Vatansever, ``Central Complex Topologies Outperform Artificial Neural Architectures in Low-Power Robotic Navigation,'' \textit{Zenodo}, DOI: 10.5281/zenodo.24375969, 2026.
\end{thebibliography}

\end{document}
